\documentclass[11pt,letterpaper]{report}
\newcommand{\Empresa}{Synaptix}
\newcommand{\EmpresaArchivo}{SYNAPTIX}
\newcommand{\RazonSocial}{Synaptix Ingeniería de Software SpA}
\newcommand{\RUTEmpresa}{76.913.482-4}
\newcommand{\CodigoLicitacion}{TFEP-01/2026}
\newcommand{\RepresentanteLegal}{Ignacio Silva}
\newcommand{\NombreOferta}{Puerto Deportivo Bahía Panitao}
\newcommand{\VersionDocumento}{0.2}
\newcommand{\FechaDocumento}{8 de octubre de 2026}
\newcommand{\ClasificacionDocumento}{Confidencial}


\usepackage{estilo/propuesta-tecnica}
% Portada vectorial reutilizable. No requiere una imagen de fondo.
\NewDocumentCommand{\PortadaPropuesta}{m m}{%
  \begin{titlepage}
    \thispagestyle{empty}
    \begin{tikzpicture}[remember picture,overlay]
      \fill[SynNavy] (current page.south west) rectangle (current page.north east);
      \fill[SynNavyLight,opacity=.52]
        ($(current page.north east)+(-4.2cm,0)$) rectangle (current page.south east);

      % Red de conexiones inspirada en la identidad digital de Synaptix.
      \draw[SynCyan,line width=2.2pt,opacity=.92]
        ($(current page.south west)+(8.2cm,2.2cm)$)
        .. controls +(3.4cm,4.0cm) and +(-4.4cm,-3.0cm) ..
        ($(current page.north east)+(-1.0cm,-4.2cm)$);
      \draw[SynBlue,line width=1.2pt,opacity=.78]
        ($(current page.south west)+(10.0cm,0.4cm)$)
        .. controls +(1.0cm,5.7cm) and +(-3.3cm,-1.0cm) ..
        ($(current page.east)+(0cm,3.0cm)$);
      \draw[SynTeal,line width=1.25pt,opacity=.70]
        ($(current page.south west)+(6.4cm,5.0cm)$)
        .. controls +(4.2cm,2.4cm) and +(-2.2cm,-4.0cm) ..
        ($(current page.north east)+(-2.0cm,-1.0cm)$);
      \draw[SynCyan!65,line width=.8pt,opacity=.62]
        ($(current page.south west)+(11.4cm,3.0cm)$)
        .. controls +(2.6cm,2.8cm) and +(-1.5cm,-3.2cm) ..
        ($(current page.north east)+(-.4cm,-7.8cm)$);

      \foreach \p/\s in {
        {($(current page.south west)+(8.2cm,2.2cm)$)}/3.2pt,
        {($(current page.south west)+(11.0cm,6.2cm)$)}/4.2pt,
        {($(current page.east)+(-1.3cm,3.0cm)$)}/3.2pt,
        {($(current page.north east)+(-2.0cm,-1.0cm)$)}/4.4pt,
        {($(current page.north east)+(-1.0cm,-4.2cm)$)}/5.0pt,
        {($(current page.north east)+(-.4cm,-7.8cm)$)}/3.1pt
      }{
        \fill[SynCyan] \p circle (\s);
        \draw[SynCyan!45,line width=.7pt] \p circle ({\s+3.2pt});
      }

      \node[anchor=north west] at ($(current page.north west)+(1.55cm,-1.35cm)$)
        {\fontsize{18}{20}\selectfont\textbf{\textcolor{SynWhite}{SYNAP}\textcolor{SynCyan}{TIX}}};
      \node[anchor=north east,text=SynCyan] at ($(current page.north east)+(-1.45cm,-1.46cm)$)
        {\fontsize{9}{11}\selectfont\bfseries INTEGRACIÓN · DATOS · DECISIONES};

      \node[anchor=west,align=left,text width=.66\paperwidth]
        at ($(current page.west)+(1.55cm,2.7cm)$) {%
          {\fontsize{13}{16}\selectfont\bfseries\textcolor{SynCyan}{PROPUESTA}}\\[2mm]
          {\fontsize{34}{38}\selectfont\bfseries\textcolor{SynWhite}{TÉCNICA}}\\[7mm]
          {\fontsize{12}{15}\selectfont\textcolor{SynWhite!82}{\NombreOferta}}
        };

      \node[anchor=west,align=left,text width=.68\paperwidth]
        at ($(current page.west)+(1.55cm,-1.0cm)$) {%
          {\fontsize{10}{12}\selectfont\bfseries\textcolor{SynCyan}{SUBDOCUMENTO #1}}\\[2mm]
          {\fontsize{19}{23}\selectfont\bfseries\textcolor{SynWhite}{#2}}
        };

      \draw[SynCyan,line width=1.5pt]
        ($(current page.south west)+(1.55cm,3.25cm)$) -- ++(3.2cm,0);
      \node[anchor=south west,align=left,text=SynWhite!74]
        at ($(current page.south west)+(1.55cm,1.35cm)$) {%
          \fontsize{9}{12}\selectfont
          \textbf{Versión} \VersionDocumento\quad
          \textbf{Fecha} \FechaDocumento\quad
          \textbf{Clasificación} \ClasificacionDocumento\\[2mm]
          \textbf{Firmante autorizado:} \FirmanteDocumento
        };
    \end{tikzpicture}
    \null
  \end{titlepage}
  \clearpage
}

\AtBeginDocument{\spanishplainpercent}
\NewDocumentCommand{\FilaDeclaracionIA}{m m m m m m}{%
  #1 & #2 & #3 & #4 & #5 & #6 \\
  \hline
}

\providecommand{\DeclaracionesIA}{%
  \FilaDeclaracionIA
    {Plantilla base}
    {OpenAI Codex}
    {Diseño y generación inicial de la plantilla LaTeX}
    {Medio}
    {Alto}
    {Equipo proponente: control humano previo a emisión; revisión pendiente.}%
  \FilaDeclaracionIA{Formato del SD8, 8 de octubre de 2026}{OpenAI Codex}{Adaptación de portada, identidad, metadatos, índice, foliación, firmas y cierre; corrección de jerarquía LaTeX en 8.1.}{Bajo}{Ninguno}{José Lara: responsable del control humano editorial previo a emisión; revisión pendiente. Alcance: formato, metadatos, índice, foliación y PDF.}
}

\newenvironment{declaracionia}{%
  \begingroup
  \setlength{\tabcolsep}{3pt}
  \begin{longtable}{L{\dimexpr.115\textwidth-2\tabcolsep\relax}L{\dimexpr.17\textwidth-2\tabcolsep\relax}L{\dimexpr.215\textwidth-2\tabcolsep\relax}L{\dimexpr.11\textwidth-2\tabcolsep\relax}L{\dimexpr.14\textwidth-2\tabcolsep\relax}L{\dimexpr.25\textwidth-2\tabcolsep\relax}}
    \caption{Declaración de uso de IA por sección y archivo asociado}
    \label{tab:declaracion-ia} \\
    \hline
    \EncabezadoTabla{Sección} &
    \EncabezadoTabla{Herramienta} &
    \EncabezadoTabla{Finalidad del uso} &
    \EncabezadoTabla{Nivel en texto} &
    \EncabezadoTabla{Nivel en diagramas} &
    \EncabezadoTabla{Revisión humana: responsable y verificación} \\
    \hline
    \endfirsthead
    \hline
    \EncabezadoTabla{Sección} &
    \EncabezadoTabla{Herramienta} &
    \EncabezadoTabla{Finalidad del uso} &
    \EncabezadoTabla{Nivel en texto} &
    \EncabezadoTabla{Nivel en diagramas} &
    \EncabezadoTabla{Revisión humana: responsable y verificación} \\
    \hline
    \endhead
}{%
  \end{longtable}
  \endgroup
}

\newcommand{\CierreSubdocumento}{%
  \clearpage
  \chapter*{Referencias}
  \addcontentsline{toc}{chapter}{Referencias}
  \markboth{Referencias}{Referencias}
  Esta lista reúne las fuentes citadas en el documento.
  \printbibliography[heading=none]
  \clearpage
  \chapter*{Declaración de uso de IA}
  \addcontentsline{toc}{chapter}{Declaración de uso de IA}
  \markboth{Declaración de uso de IA}{Declaración de uso de IA}
  Synaptix registra por sección las herramientas de IA generativa utilizadas, su finalidad, el grado de intervención y la revisión humana. Ninguno indica ausencia de uso. En texto, Bajo corresponde a corrección de redacción; Medio, a borradores reescritos y verificados por el equipo; Alto, a generación sustancial. En diagramas, Bajo corresponde a sugerencias de formato; Medio, a generación asistida desde un modelo del equipo; Alto, a generación sustancial. El registro se consolidará en el Formulario A-6. Synaptix asume la responsabilidad por la exactitud y coherencia de la propuesta.

  La Tabla~\ref{tab:declaracion-ia} registra el uso declarado por sección y formulario asociado.
  \begin{declaracionia}
    \DeclaracionesIA
  \end{declaracionia}
  \FuenteTabla{Registro de uso de IA de este subdocumento.}
}

% Convenciones comunes del SD8 y T-16.
\setlength{\headheight}{23pt}
\newcommand{\ArchivoTDieciseis}{\texttt{SYNAPTIX-Formulario-T-16.pdf}}
\newcommand{\ArchivoSDOcho}{\texttt{SYNAPTIX-Subdocumento8.pdf}}
\AtBeginBibliography{\RaggedRight\setlength{\emergencystretch}{3em}}
\NewDocumentCommand{\TituloFormularioRiesgos}{m}{\chapter*{#1}\addcontentsline{toc}{chapter}{#1}\markboth{#1}{#1}}
\RenewDocumentCommand{\FuenteTabla}{m}{\FuenteFigura{#1}\par\smallskip}

\addbibresource{config/riesgos.bib}
\newcommand{\FormularioRiesgos}{T-16}
% Nuevas intervenciones; preserva declaraciones históricas del SD8.
\newcommand{\ControlHumanoRiesgos}{José Lara: responsable del control humano técnico y editorial previo a emisión; revisión pendiente. Alcance: fuentes, criterios, riesgos, cálculos, concordancia, respuestas y PDF.}
\ifdefined\FormularioRiesgos
 \renewcommand{\DeclaracionesIA}{%
  \FilaDeclaracionIA{T-16, 2026-10-08}{OpenAI Codex}{Formulación de dieciséis riesgos, matriz oficial, fichas, objetivos residuales y trazabilidad a SD7.}{Alto}{Ninguno}{\ControlHumanoRiesgos}
 }
\else
 \let\DeclaracionesAnterioresRiesgos\DeclaracionesIA
 \renewcommand{\DeclaracionesIA}{%
  \DeclaracionesAnterioresRiesgos
  \FilaDeclaracionIA{8.1, 2026-10-08}{OpenAI Codex}{Desarrollo del plan ISO 31000, roles, escalas, registro, seguimiento y figura del proceso.}{Alto}{Alto}{\ControlHumanoRiesgos}
  \FilaDeclaracionIA{8.2, 2026-10-08}{OpenAI Codex}{RBS, priorización, modos de falla y sensibilidades de plazo/esfuerzo con supuestos explícitos.}{Alto}{Ninguno}{\ControlHumanoRiesgos}
  \FilaDeclaracionIA{8.3, 2026-10-08}{OpenAI Codex}{Respuestas, costo-beneficio sin precios, contingencia, gestión, eficacia y cierre.}{Alto}{Ninguno}{\ControlHumanoRiesgos}
  \FilaDeclaracionIA{T-16, 2026-10-08}{OpenAI Codex}{Registro y formulario independientes; conservación de responsables PL-R y distinción de residual objetivo/observado.}{Alto}{Ninguno}{\ControlHumanoRiesgos}
 }
\fi

\renewcommand{\thesection}{T-16.\arabic{section}}
\renewcommand{\thetable}{T-16.\arabic{table}}
\setlength{\cftsecnumwidth}{4.5em}
\setlength{\cftsubsecnumwidth}{6em}
\begin{document}
\pagenumbering{arabic}
\PortadaFormulario{T-16}{PLAN DE RIESGOS}
\tableofcontents\clearpage
\TituloFormularioRiesgos{Plan de riesgos}
Synaptix presenta la matriz y las respuestas de \ArchivoSDOcho{}, con los ocho campos oficiales de T-16 y trazabilidad adicional a SD7. Las valoraciones son hipótesis de oferta y los controles se verificarán en ejecución; no constituyen probabilidades medidas ni pruebas realizadas. \parencites[Formulario T-16, p.~63]{sd8-ba}[RT-19.04, p.~33]{sd8-bt}[cap.~10, p.~26; cap.~13, pp.~28--30; cap.~15, pp.~33--36]{sd8-c07}
\section{Matriz oficial y criterios}
La Tabla~\ref{tab:t16-matriz} utiliza $P/I$ ordinales de uno a cinco y $E=PI$. Los intervalos bajo 1--4, moderado 5--9, apreciable 10--14, alto 15--19 y crítico 20--25 priorizan el tratamiento. Impacto 5 en personas, integridad o privacidad exige revisión competente aunque $E$ sea menor a 15; ninguna aceptación residual permite incumplir una cláusula. Los objetivos residuales de las fichas son metas sujetas a evidencia, no resultados observados.
% Fuente: registro_riesgos.csv; valores iniciales y tratamientos de oferta.
\begin{landscape}
\begingroup\fontsize{9}{11}\selectfont\setlength{\tabcolsep}{2pt}
\begin{longtable}{L{\dimexpr.035\linewidth-2\tabcolsep\relax}L{\dimexpr.215\linewidth-2\tabcolsep\relax}L{\dimexpr.13\linewidth-2\tabcolsep\relax}L{\dimexpr.06\linewidth-2\tabcolsep\relax}L{\dimexpr.06\linewidth-2\tabcolsep\relax}L{\dimexpr.06\linewidth-2\tabcolsep\relax}L{\dimexpr.265\linewidth-2\tabcolsep\relax}L{\dimexpr.175\linewidth-2\tabcolsep\relax}}
\caption{Plan de riesgos: campos oficiales del Formulario T-16}\label{tab:t16-matriz}\\\hline
\EncabezadoTabla{N°} & \EncabezadoTabla{Riesgo} & \EncabezadoTabla{Categoría} & \EncabezadoTabla{Prob.} & \EncabezadoTabla{{\fontsize{8}{9}\selectfont Impacto}} & \EncabezadoTabla{Expos.} & \EncabezadoTabla{Estrategia / Mitigación} & \EncabezadoTabla{Responsable}\\\hline\endfirsthead
\hline\EncabezadoTabla{N°} & \EncabezadoTabla{Riesgo} & \EncabezadoTabla{Categoría} & \EncabezadoTabla{Prob.} & \EncabezadoTabla{{\fontsize{8}{9}\selectfont Impacto}} & \EncabezadoTabla{Expos.} & \EncabezadoTabla{Estrategia / Mitigación} & \EncabezadoTabla{Responsable}\\\hline\endhead
1 & R-01: Cancelación de trabajo programado por marejadas & Proyecto & 4 & 5 & 20 & Preparar ventanas alternativas y revisar la red completa con suspensión de todos los frentes; proteger el margen anterior al hito. & José Lara Arce (Implantación)\\\hline
2 & R-02: Pérdida conjunta de enlace, energía o nodo local & Operación & 3 & 5 & 15 & Comprobar caminos de energía, exclusión de escrituras, respaldo y autonomías 24/12/8 horas con mezcla operacional. & José Lara Arce (SRE)\\\hline
3 & R-03: Conflicto de autoridad o pérdida en sincronización & Técnico & 3 & 5 & 15 & Probar idempotencia, autoridad por período, conflictos y reconexión tras veinticuatro horas; conservar originales y versiones. & Fernando Guerrero (Arquitectura)\\\hline
4 & R-04: Contraparte no disponible para decisión o acta & Organizacional & 4 & 5 & 20 & Acordar referentes y suplentes por proceso, revisión continua y agenda de decisiones antes de la ventana; entregar expedientes completos. & Ignacio Silva (Proyecto)\\\hline
5 & R-05: Capacidad compartida y factores UCP deteriorados & Proyecto & 3 & 4 & 12 & Nivelar por persona y día, reservar revisión y verificar competencias, disponibilidad y puestos de refuerzo. & Ignacio Silva (Proyecto)\\\hline
6 & R-06: Obsolescencia, variante o suministro incompatibles & Técnico & 3 & 4 & 12 & Verificar variante, soporte, licencias, interfaces y sustitución antes de compra; mantener inventario y plan de actualización. & Fernando Guerrero (Arquitectura)\\\hline
7 & R-07: Bloqueo por proveedor en la salida & Técnico & 3 & 4 & 12 & Versionar contratos, IaC y exportaciones; ensayar alternativas, recifrado y límites de portabilidad; actualizar plan de reversibilidad. & Fernando Guerrero (Arquitectura)\\\hline
8 & R-08: Saturación en peak o crecimiento de telemetría & Técnico & 3 & 4 & 12 & Probar mezcla, carga/estrés, degradación y recuperación; proyectar capacidad trimestral y límites de telemetría. & Fernando Guerrero (Arquitectura)\\\hline
9 & R-09: Compromiso de dependencias o entrega & Seguridad & 3 & 5 & 15 & Ramas protegidas, SAST/SCA/DAST, análisis de imágenes, SBOM y aislamiento de firma SLSA L3; intrusión independiente. & Bruno Toro (Seguridad)\\\hline
10 & R-10: Exposición de menores, salud o datos protegidos & Seguridad & 3 & 5 & 15 & Minimización, cifrado de campo, permisos por finalidad y segregación; verificar rechazo, auditoría y datos de prueba protegidos. & Bruno Toro (Seguridad)\\\hline
11 & R-11: Intervención o equipo que expone personas & Operación & 3 & 5 & 15 & Recepción de equipos y montaje marino, prueba con guantes/humedad y segregación; respetar obras/compras del CLIENTE y ventanas. & Fernando Guerrero (Arquitectura)\\\hline
12 & R-12: Migración o retorno no conciliados & Proyecto & 3 & 5 & 15 & Perfilar fuentes con referente, conservar originales y realizar dos ensayos completos por etapa con conciliación y retorno. & Davor Serey (Calidad)\\\hline
13 & R-13: Cobertura insuficiente del servicio gestionado & Operación & 3 & 4 & 12 & Acreditar puestos y suplencia NOC/SOC/mesa bilingüe, medir colas y comprobar escalamiento 24 por 7 para emergencias. & José Lara Arce (Operación)\\\hline
14 & R-14: Rotación estacional y adopción incompleta & Organizacional & 4 & 4 & 16 & Formación por perfil/turno en noviembre, ejercicios de tarea, suplentes y mentoría; registrar evaluación individual. & José Lara Arce (Implantación)\\\hline
15 & R-15: Pilotos o evidencia ambiental fuera de ventana & Proyecto & 3 & 4 & 12 & Mantener fechas y dependencias de SD13, habilitar medición temprana y verificar datos/pilotos con criterios propios. & Ignacio Silva (Proyecto)\\\hline
16 & R-16: Fuentes, licencias o custodia incompletas & Proyecto & 3 & 4 & 12 & Repositorio del CLIENTE, verificación de depósito en cada liberación y contrato de custodia independiente con condiciones de liberación. & Ignacio Silva (Proyecto)\\\hline
17 & R-IN1-01: Actualización insuficiente por contratistas & Organizacional & 4 & 4 & 16 & Formar contratistas y marina; comprobar tareas y recordatorios con fuente de cada cambio. & José Lara Arce (Implantación)\\\hline
18 & R-IN1-02: Participación insuficiente de armadores & Organizacional & 3 & 3 & 9 & Orientar e ingresar solicitudes de forma asistida sin crear demanda artificial; conservar excluidas y no respondidas. & José Lara Arce (Implantación)\\\hline
19 & R-IN1-03: Integración incompatible de solicitudes y órdenes & Técnico & 3 & 4 & 12 & Probar contratos, idempotencia, concurrencia, habilitación y recuperación de avisos; validar vínculo con orden autorizada. & Fernando Guerrero (Arquitectura)\\\hline
20 & R-IN1-04: Acceso a solicitudes de otra embarcación & Seguridad & 2 & 5 & 10 & Validar vínculo vigente y ámbito en servidor; ensayar separación entre armadores y contratistas con auditoría. & Bruno Toro (Seguridad)\\\hline
21 & R-IN2-01: Cohorte de invernada insuficiente & Organizacional & 3 & 4 & 12 & Convocar y formar armadores; confirmar planes y disponibilidad sin imponer contratista ni trabajos. & José Lara Arce (Implantación)\\\hline
22 & R-IN2-02: Atraso del contratista respecto de botadura & Proyecto & 4 & 4 & 16 & Coordinar líder funcional y administración de varadero; confirmar disponibilidad y los cinco lotes de cuatro embarcaciones. & José Lara Arce (Implantación)\\\hline
23 & R-IN2-03: Suspensión meteorológica durante invernada & Proyecto & 4 & 5 & 20 & Integrar ventana y aviso con R-01; comprobar reserva por campaña y red de todos los frentes antes de seleccionar trabajo. & Ignacio Silva (Proyecto)\\\hline
24 & R-IN2-04: IN1 o interfaces no disponibles para IN2 & Técnico & 3 & 4 & 12 & Verificar contrato, estados y disponibilidad de IN1 antes de habilitar el plan preventivo. & Fernando Guerrero (Arquitectura)\\\hline
25 & R-IN3-01: Campos OCR incorrectos por formato o calidad & Técnico & 4 & 4 & 16 & Líder de Datos aplica reglas, fragmentos de origen y evaluación por formato; conservar original y propuesta separada del dato aprobado. & Fernando Guerrero (Arquitectura; coordinación con Líder de Datos)\\\hline
26 & R-IN3-02: Revisión humana insuficiente de extracción & Organizacional & 3 & 4 & 12 & Capacitar revisores con casos erróneos; exigir confirmación de campos y comprobar registros aprobados por QA. & José Lara Arce (Implantación; con verificación de Calidad)\\\hline
27 & R-IN3-03: Tratamiento documental no autorizado & Seguridad & 3 & 5 & 15 & Comprobar ubicación, retención, permisos y condiciones de procesamiento antes de enviar; conservar trazabilidad y acceso mínimo. & Bruno Toro (Seguridad)\\\hline
28 & R-IN3-04: Extractor o conectividad interrumpidos & Operación & 3 & 3 & 9 & Supervisar fallos y ensayar alternativa manual, reintentos e idempotencia; separar pendientes de antecedentes aprobados. & José Lara Arce (Operación/SRE)\\\hline
29 & R-IN3-05: Beneficio de extracción insuficiente & Proyecto & 3 & 3 & 9 & Medir manual/asistido con revisión incluida, por formato; registrar esfuerzo y contrastar costos incrementales exclusivamente en OE. & Ignacio Silva (Proyecto)\\\hline
30 & R-IN4-01: Atribución inválida o pago por comparación no equivalente & Proyecto & 3 & 5 & 15 & Fijar ficha y composición antes de medir; conservar intervalos activos, fuentes, exclusiones y cálculo; separar preparación Synaptix de verificación cliente. & Ignacio Silva (Proyecto)\\\hline
31 & R-IN4-02: Omisión de casos difíciles o cierre apresurado por incentivo & Proyecto & 3 & 5 & 15 & Calidad contrasta universo elegible con AUT/VAR/CTR, incluidas abiertas y excluidas; fija complejidad y plazo antes de resultados; compara calidad por condición. & Ignacio Silva (Proyecto; con verificación de Calidad)\\\hline
32 & R-IN5-01: Reducción ambiental aparente por datos o selección sesgados & Técnico & 3 & 4 & 12 & Comparar mismas embarcaciones, conservar asociación histórica de medidor y permanencia validada; fijar criterios y registrar faenas, clima y resultados adversos. & Ignacio Silva (Proyecto; con verificación de Calidad)\\\hline
33 & R-IN5-02: Participación o continuidad insuficientes en planes ambientales & Organizacional & 3 & 3 & 9 & Validar acciones breves con armadores y charter; explicar responsabilidades y retiro sin pérdida de servicio básico. & José Lara Arce (Implantación)\\\hline
34 & R-IN5-03: Acceso entre flotas o uso fuera de finalidad ambiental & Seguridad & 2 & 5 & 10 & Verificar vínculo vigente, finalidad y autorización por embarcación; ensayar flotas separadas y propagación de restricciones con auditoría. & Bruno Toro (Seguridad)\\\hline
\end{longtable}
\FuenteTabla{Registro de Synaptix; campos del T-16, Bases Administrativas, p.~63; escalas de SD8.}
\endgroup
\end{landscape}

R-01 y R-04 requieren decisión temprana por su exposición 20. Los controles de los demás riesgos conservan sus puertas obligatorias. Un consumo de reserva se registra una sola vez aunque el evento afecte varias fichas.
\section{Respuestas, disparadores y trazabilidad}
Las fichas siguientes completan causa, contingencia, plazo, evidencia, paquetes y reserva de la matriz oficial. El objetivo residual depende de eficacia por comprobar; los responsables de los riesgos PL-R de SD7 permanecen en su registro de origen. Cada acción se valida en capacidad y cuenta de costo antes de ejecutarse.

\subsection{R-01: Cancelación de trabajo programado por marejadas}
Avisos de la autoridad marítima y ventanas costeras limitadas. Se suspende toda actividad programada, incluida construcción aislada o remota. Se consumen reservas y puede faltar capacidad antes de H3 o del corte de etapa. La evaluación inicial es $P=4$, $I=5$ y $E=20$; el objetivo posterior al tratamiento es $P_o=4$, $I_o=3$ y $E_o=12$, sin acreditación de residual observado.

José Lara Arce (Implantación) responde por el riesgo. Disparador: Aviso oficial de marejadas; afectación de cualquier frente programado. Preparar ventanas alternativas y revisar la red completa con suspensión de todos los frentes; proteger el margen anterior al hito. Contingencia: Suspender de inmediato, registrar capacidad afectada y aportar refuerzo propio donde sea admisible; reanudar sólo tras cese y autorización. Plazo: Preparación antes de cada campaña; suspensión inmediata; evaluación de impacto el mismo día del aviso.

Se exige la evidencia siguiente: Aviso, bloqueo, red actualizada, capacidad, reserva consumida y autorización de reanudación. Paquetes vinculados: PL-R01; 1.2.3; 1.7.2.5; ventanas y cortes de SD7. Reserva: Diez días hábiles por campaña marina, una sola vez; no cubren automáticamente otros frentes. Fundamento: C07 cap. 10, restricción 12, p. 26; numeral 13.2, p. 29; RT-10.05, p. 34; SD7 PL-R01. Estado: Identificado; tratamiento comprometido; residual objetivo no verificado. \parencite[cap.~10, restricción 12, p.~26; numeral 13.2, p.~29; RT-10.05, p.~34]{sd8-c07}

\subsection{R-02: Pérdida conjunta de enlace, energía o nodo local}
Dependencias físicas y fallos correlacionados en el recinto. La plataforma local deja de sostener los flujos críticos durante un aislamiento. Se comprometen continuidad, registro e integridad de salidas, escuela o combustible. La evaluación inicial es $P=3$, $I=5$ y $E=15$; el objetivo posterior al tratamiento es $P_o=2$, $I_o=5$ y $E_o=10$, sin acreditación de residual observado.

José Lara Arce (SRE) responde por el riesgo. Disparador: Falla de conmutación, autonomía o recuperación en ensayo; incidente crítico. Comprobar caminos de energía, exclusión de escrituras, respaldo y autonomías 24/12/8 horas con mezcla operacional. Contingencia: Activar recuperación e incidente, proteger originales y usar respaldo operacional autorizado; impedir aceptación si no se alcanza el compromiso. Plazo: Antes de cada producción; DR/resiliencia al menos semestral; activación inmediata ante incidente.

Se exige la evidencia siguiente: Ensayo de autonomía, fallo y DR real, tiempos, último dato y conciliación. Paquetes vinculados: 1.2; 1.3; 1.10.2.e.5/6/7; 1.12.4. Reserva: Usar respuesta de continuidad ya programada; refuerzo adicional se dimensiona sin doble imputación. Fundamento: C07 RT-03.10/13, p. 33; BT cap. 20, pp. 34–35; SD4 continuidad. Estado: Identificado; tratamiento comprometido; residual objetivo no verificado. \parencites[RT-03.10/13, p.~33]{sd8-c07}[cap.~20, pp.~34--35]{sd8-bt}

\subsection{R-03: Conflicto de autoridad o pérdida en sincronización}
Reintentos, doble escritura o reglas cloud/local incompatibles. Un hecho se duplica, pierde o interpreta con autoridad incorrecta al reconectar. Se altera presencia, amarra, retiro, combustible o saldo. La evaluación inicial es $P=3$, $I=5$ y $E=15$; el objetivo posterior al tratamiento es $P_o=2$, $I_o=5$ y $E_o=10$, sin acreditación de residual observado.

Fernando Guerrero (Arquitectura) responde por el riesgo. Disparador: Conciliación distinta de cero, duplicado, permiso inválido o sincronización mayor a treinta minutos. Probar idempotencia, autoridad por período, conflictos y reconexión tras veinticuatro horas; conservar originales y versiones. Contingencia: Bloquear la promoción y escrituras incompatibles, preservar el diario y corregir con conciliación; repetir evidencia invalidada. Plazo: Antes de aceptar contratos y cada producción; revisión ante cambio de esquema o autoridad.

Se exige la evidencia siguiente: Casos negativos, diario, conciliación y tiempos por digest. Paquetes vinculados: 1.3; 1.4 a 1.6; 1.10.2.e.5; contratos SD5. Reserva: Corrección en paquete causante; escenario PL-R02 no aporta HH adicionales por sí solo. Fundamento: BT RT-04.10, p. 11; C07 RT-03.13, p. 33; SD4/SD5 autoridad. Estado: Identificado; tratamiento comprometido; residual objetivo no verificado. \parencites[RT-04.10, p.~11]{sd8-bt}[RT-03.13, p.~33]{sd8-c07}

\subsection{R-04: Contraparte no disponible para decisión o acta}
Cliente sin TI, turnos operacionales y validación final sin holgura. No se resuelve una decisión o no se firma el cierre con evidencia completa. Se bloquea aceptación o corte, especialmente al inicio del mes 21. La evaluación inicial es $P=4$, $I=5$ y $E=20$; el objetivo posterior al tratamiento es $P_o=3$, $I_o=4$ y $E_o=12$, sin acreditación de residual observado.

Ignacio Silva (Proyecto) responde por el riesgo. Disparador: Referente/suplente sin disponibilidad; observación impeditiva o ausencia de firma programada. Acordar referentes y suplentes por proceso, revisión continua y agenda de decisiones antes de la ventana; entregar expedientes completos. Contingencia: Escalar, mantener legado y bloquear corte sin acta; aportar subsanación propia y registrar atraso sin presumir aceptación tácita. Plazo: Referentes en arranque; agenda antes de cada recepción; aviso de desvío según RT-19.10.

Se exige la evidencia siguiente: Agenda y participación confirmadas, entrega, observaciones, decisiones y acta expresa. Paquetes vinculados: PL-R04; 1.1.1.1; 1.1.4; 1.11.5; ACTA-MB-E2. Reserva: Cero holgura en acta/corte del primer día de 21; no reducir derechos de revisión. Fundamento: BA arts. 17/18, pp. 12–13, y 71, pp. 37–38; C07 13.3, pp. 29–30; SD7 PL-R04. Estado: Identificado; tratamiento comprometido; residual objetivo no verificado. \parencites[arts.~17/18, pp.~12--13; art.~71, pp.~37--38]{sd8-ba}[numeral 13.3, pp.~29--30]{sd8-c07}

\subsection{R-05: Capacidad compartida y factores UCP deteriorados}
Funciones coincidentes, refuerzos no habilitados o factores de ambiente distintos del modelo. El esfuerzo o la revisión supera la capacidad efectiva de T-15. Se forman colas y quedan paquetes sin ejecución o revisión antes de sus puertas. La evaluación inicial es $P=3$, $I=4$ y $E=12$; el objetivo posterior al tratamiento es $P_o=2$, $I_o=4$ y $E_o=8$, sin acreditación de residual observado.

Ignacio Silva (Proyecto) responde por el riesgo. Disparador: Agenda sobreasignada, revisor ausente o escenario funcional que pasa de CF20 a CF28. Nivelar por persona y día, reservar revisión y verificar competencias, disponibilidad y puestos de refuerzo. Contingencia: Recalcular esfuerzo y red, limitar WIP y financiar capacidad propia adicional sin reducir pruebas ni mover hitos. Plazo: Cada selección quincenal y antes de habilitar un frente o refuerzo.

Se exige la evidencia siguiente: Asignación sin sobrecarga, comparación UCP y provisión por persona. Paquetes vinculados: PL-R05; 1.1.2.2; 1.4 a 1.6; T-15. Reserva: 12.708 HH base y 19.640,50 HH de sensibilidad son escenarios distintos, no bolsas simultáneas. Fundamento: SD7 PL-R05, asignación y sensibilidad UCP; SD6 WIP. Estado: Identificado; tratamiento comprometido; residual objetivo no verificado.

\subsection{R-06: Obsolescencia, variante o suministro incompatibles}
Fin de soporte, ficha incorrecta o componente no apto para el entorno marino. La variante recibida no sostiene compatibilidad, protección o soporte contractual. Se requiere sustitución y puede demorarse H3 o una reposición. La evaluación inicial es $P=3$, $I=4$ y $E=12$; el objetivo posterior al tratamiento es $P_o=2$, $I_o=3$ y $E_o=6$, sin acreditación de residual observado.

Fernando Guerrero (Arquitectura) responde por el riesgo. Disparador: Variante no trazable, fin de soporte dentro del uso o ensayo físico no conforme. Verificar variante, soporte, licencias, interfaces y sustitución antes de compra; mantener inventario y plan de actualización. Contingencia: Rechazar componente, seleccionar sustitución equivalente o superior y repetir integración/recepción; coordinar adquisición del CLIENTE. Plazo: Antes de cada compra y H3; revisión trimestral de soporte y ante aviso del fabricante.

Se exige la evidencia siguiente: Ficha y variante, soporte, compatibilidad y aceptación del sustituto. Paquetes vinculados: 1.2; 1.7.2; EXT-EQUIPOS/LICENCIAS; T-11. Reserva: Margen de abastecimiento de SD7; reemplazo se dimensiona sin atribuir catálogo a una hipótesis. Fundamento: BT 1.6, p. 4, y cap. 8, pp. 18–20; C07 caps. 10/11, pp. 26–27; SD4 inventario. Estado: Identificado; tratamiento comprometido; residual objetivo no verificado. \parencites[numeral 1.6, p.~4; cap.~8, pp.~18--20]{sd8-bt}[caps.~10/11, pp.~26--27]{sd8-c07}

\subsection{R-07: Bloqueo por proveedor en la salida}
Servicios administrados, formatos o claves no exportables. No se puede transferir datos, fuentes, identidad o configuración con continuidad. Se prolonga dependencia o falla la transferencia contractual. La evaluación inicial es $P=3$, $I=4$ y $E=12$; el objetivo posterior al tratamiento es $P_o=2$, $I_o=3$ y $E_o=6$, sin acreditación de residual observado.

Fernando Guerrero (Arquitectura) responde por el riesgo. Disparador: Cambio de servicio/licencia, exportación no íntegra o incompatibilidad del destino. Versionar contratos, IaC y exportaciones; ensayar alternativas, recifrado y límites de portabilidad; actualizar plan de reversibilidad. Contingencia: Mantener lectura del origen, adaptar servicios no portables y ejecutar transferencia autorizada con comprobación y retorno. Plazo: Plan inicial dentro de noventa días; actualización anual; salida acompañada en meses 53 a 56.

Se exige la evidencia siguiente: Inventario, exportación conciliada, procedimientos y recepción de accesos. Paquetes vinculados: 1.1.3; 1.12.7/8; SD4 gobierno\_nube\_rt. Reserva: La sustitución de plataforma estimada en SD4 no se confunde con transferencia ya presupuestada. Fundamento: BA art. 84, p. 43; SD4/SD7 reversibilidad y salida. Estado: Identificado; tratamiento comprometido; residual objetivo no verificado. \parencite[art.~84, p.~43]{sd8-ba}

\subsection{R-08: Saturación en peak o crecimiento de telemetría}
Concentración operacional y frecuencia de medición mayor a la dimensionada. Se incumplen tiempos o aparecen colas sin recuperación controlada. Se demora un flujo crítico o se pierde visibilidad de escuela y embarcaciones. La evaluación inicial es $P=3$, $I=4$ y $E=12$; el objetivo posterior al tratamiento es $P_o=2$, $I_o=3$ y $E_o=6$, sin acreditación de residual observado.

Fernando Guerrero (Arquitectura) responde por el riesgo. Disparador: Carga a 1,5 veces peak incumple umbral; alerta de capacidad o cambio de frecuencia. Probar mezcla, carga/estrés, degradación y recuperación; proyectar capacidad trimestral y límites de telemetría. Contingencia: Aplicar límites y escalamiento autorizados, dimensionar recursos y repetir pruebas antes de promover. Plazo: Antes de cada producción y proyección trimestral en Operación.

Se exige la evidencia siguiente: Curvas, saturación, recursos, tiempos de negocio y recuperación. Paquetes vinculados: 1.10.2.e.4; 1.12.3.5; SD4 capacidad. Reserva: Capacidad de prueba y ajuste del paquete; no usar promedios como peak. Fundamento: BT cap. 9, p. 21, y cap. 20, pp. 34–35; C07 14.1/14.2, pp. 31–32, y RT-09.01/02, p. 34. Estado: Identificado; tratamiento comprometido; residual objetivo no verificado. \parencites[cap.~9, p.~21; cap.~20, pp.~34--35]{sd8-bt}[numerales 14.1/14.2, pp.~31--32; RT-09.01/02, p.~34]{sd8-c07}

\subsection{R-09: Compromiso de dependencias o entrega}
Origen no verificado, vulnerabilidad, secreto expuesto o procedencia falsificable. Se incorpora o promueve un artefacto inseguro. Se afecta confidencialidad, integridad o continuidad de un servicio. La evaluación inicial es $P=3$, $I=5$ y $E=15$; el objetivo posterior al tratamiento es $P_o=2$, $I_o=5$ y $E_o=10$, sin acreditación de residual observado.

Bruno Toro (Seguridad) responde por el riesgo. Disparador: Hallazgo crítico/alto, firma o SBOM inválida, secreto o prueba incompleta. Ramas protegidas, SAST/SCA/DAST, análisis de imágenes, SBOM y aislamiento de firma SLSA L3; intrusión independiente. Contingencia: Bloquear promoción, revocar/rotar credenciales, aislar artefacto y corregir; respetar máximos de siete/quince/treinta días. Plazo: Cada commit/liberación; intrusión antes de cada producción y anual.

Se exige la evidencia siguiente: Origen, análisis, firma/procedencia, informe íntegro y reexamen. Paquetes vinculados: 1.3; 1.10.2.e.8; 1.12.5; PRUEBAS-6. Reserva: Corrección y reensayo por paquete; no aceptar una falla crítica por consumir reserva. Fundamento: BA 4.3, pp. 5–6; BT RT-11.04, p. 23, y RT-11.20/22–27, p. 24; SD4 SEC-1.0. Estado: Identificado; tratamiento comprometido; residual objetivo no verificado. \parencites[art.~4.3, pp.~5--6]{sd8-ba}[RT-11.04, p.~23; RT-11.20/22--27, p.~24]{sd8-bt}

\subsection{R-10: Exposición de menores, salud o datos protegidos}
Permisos amplios, registros inseguros o datos reales en pruebas. Una persona consulta o divulga datos que no requiere para su función. Se comprometen privacidad, seguridad de menores y obligaciones de tratamiento. La evaluación inicial es $P=3$, $I=5$ y $E=15$; el objetivo posterior al tratamiento es $P_o=2$, $I_o=5$ y $E_o=10$, sin acreditación de residual observado.

Bruno Toro (Seguridad) responde por el riesgo. Disparador: Prueba de autorización fallida, exportación indebida o incidente de datos. Minimización, cifrado de campo, permisos por finalidad y segregación; verificar rechazo, auditoría y datos de prueba protegidos. Contingencia: Contener acceso, preservar evidencia, activar incidente y notificaciones aplicables y corregir antes de rehabilitar. Plazo: Antes de habilitar perfiles y cada producción; contención inmediata ante incidente.

Se exige la evidencia siguiente: Casos de acceso positivo/negativo, auditoría y revisión de exportaciones. Paquetes vinculados: 1.3; contratos SD5; 1.10.2.e.8; 1.12.3.3. Reserva: Respuesta de seguridad incluida; no se acepta excepción a protección obligatoria. Fundamento: C07 cap. 10, restricciones 3/14, p. 26, y RT-11.10, p. 34; BT cap. 11, pp. 23–24, y cap. 16, pp. 29–31. Estado: Identificado; tratamiento comprometido; residual objetivo no verificado. \parencites[cap.~10, restricciones 3/14, p.~26; RT-11.10, p.~34]{sd8-c07}[cap.~11, pp.~23--24; cap.~16, pp.~29--31]{sd8-bt}

\subsection{R-11: Intervención o equipo que expone personas}
Flexión, agua, combustible clasificado o interacción durante maniobra. Una instalación o interfaz resulta insegura en pantalán, surtidor o varadero. Se expone a personas junto a cabos, cargas o menores. La evaluación inicial es $P=3$, $I=5$ y $E=15$; el objetivo posterior al tratamiento es $P_o=2$, $I_o=5$ y $E_o=10$, sin acreditación de residual observado.

Fernando Guerrero (Arquitectura) responde por el riesgo. Disparador: Comportamiento físico no acreditado, acceso inseguro o interacción exigida durante maniobra. Recepción de equipos y montaje marino, prueba con guantes/humedad y segregación; respetar obras/compras del CLIENTE y ventanas. Contingencia: Suspender intervención, retirar componente inseguro del servicio y seleccionar alternativa verificada; no reducir protección. Plazo: Antes de instalación y habilitación; suspensión inmediata ante condición insegura.

Se exige la evidencia siguiente: Ficha de variante, montaje, prueba segura y autorización de acceso. Paquetes vinculados: 1.2.3; 1.7.2; EXT-RECINTO/EQUIPOS; formación. Reserva: Ventana alternativa y reemplazo dimensionados; ninguna reserva habilita maniobra insegura. Fundamento: C07 cap. 10, restricciones 1/9/10, p. 26; RT-03.24, p. 33; RT-13.08, p. 35; SD4 instalación. Estado: Identificado; tratamiento comprometido; residual objetivo no verificado. \parencite[cap.~10, restricciones 1/9/10, p.~26; RT-03.24, p.~33; RT-13.08, p.~35]{sd8-c07}

\subsection{R-12: Migración o retorno no conciliados}
Legado sin soporte/documentación y conocimiento concentrado. Los ensayos encuentran diferencias, errores de transformación o retorno no íntegro. Se repiten pruebas y puede faltar tiempo de revisión antes del corte. La evaluación inicial es $P=3$, $I=5$ y $E=15$; el objetivo posterior al tratamiento es $P_o=2$, $I_o=4$ y $E_o=8$, sin acreditación de residual observado.

Davor Serey (Calidad) responde por el riesgo. Disparador: Diferencia no explicada, volumen incompleto, tiempo excedido o prueba de retorno fallida. Perfilar fuentes con referente, conservar originales y realizar dos ensayos completos por etapa con conciliación y retorno. Contingencia: Bloquear corte, corregir con Datos/Arquitectura y repetir evidencia invalidada sin usar aceptación tácita. Plazo: Dos ensayos antes de cada migración definitiva; revisión y subsanación conforme BA 18.

Se exige la evidencia siguiente: Recuentos, diferencias resueltas, tiempos y retorno probado por versión. Paquetes vinculados: PL-R02/03; 1.10.1; ENS-VENT-E1/E2; datos SD5. Reserva: Escenario de seis días adicionales (10,10,16); no cubre automáticamente repetición completa ni HH de corrección. Fundamento: BA arts. 17/18, pp. 12–13; BT 20.1, pp. 34–35; C07 13.3, pp. 29–30; SD7 PL-R02/03. Estado: Identificado; tratamiento comprometido; residual objetivo no verificado. \parencites[arts.~17/18, pp.~12--13]{sd8-ba}[sección 20.1, pp.~34--35]{sd8-bt}[numeral 13.3, pp.~29--30]{sd8-c07}

\subsection{R-13: Cobertura insuficiente del servicio gestionado}
Turnos, suplencias, inglés o especialistas sin provisión efectiva. La atención no sostiene horario, emergencia o SLA de temporada. Se prolonga incidente y el CLIENTE sin TI queda sin apoyo competente. La evaluación inicial es $P=3$, $I=4$ y $E=12$; el objetivo posterior al tratamiento es $P_o=2$, $I_o=4$ y $E_o=8$, sin acreditación de residual observado.

José Lara Arce (Operación) responde por el riesgo. Disparador: Turno sin relevo, proveedor no habilitado o indicador de atención fuera del compromiso. Acreditar puestos y suplencia NOC/SOC/mesa bilingüe, medir colas y comprobar escalamiento 24 por 7 para emergencias. Contingencia: Activar suplente y especialista, ampliar cobertura propia y corregir dimensionamiento sin recargo por incumplimiento. Plazo: Antes de servicio; revisión mensual y preparación previa a cada temporada.

Se exige la evidencia siguiente: Cuadrantes, contratos/habilitación efectivos, métricas y ejercicios de escalamiento. Paquetes vinculados: 1.12.1; 1.12.3; T-15 cobertura externa. Reserva: Capacidad externa reservada separada de HH de los titulares; suplencia no se presume contratada. Fundamento: C07 cap. 10, restricción 7, p. 26, y RT-21.06, p. 36; SD7 cobertura. Estado: Identificado; tratamiento comprometido; residual objetivo no verificado. \parencite[cap.~10, restricción 7, p.~26; RT-21.06, p.~36]{sd8-c07}

\subsection{R-14: Rotación estacional y adopción incompleta}
Personal nuevo y monitores renovados, con capacitación restringida en peak. Una persona llega a la ola sin competencia o certificación requerida. Se usa incorrectamente el flujo o se bloquea la habilitación y cierre. La evaluación inicial es $P=4$, $I=4$ y $E=16$; el objetivo posterior al tratamiento es $P_o=3$, $I_o=3$ y $E_o=9$, sin acreditación de residual observado.

José Lara Arce (Implantación) responde por el riesgo. Disparador: Cohorte incompleta, evaluación no aprobada o alta posterior a formación. Formación por perfil/turno en noviembre, ejercicios de tarea, suplentes y mentoría; registrar evaluación individual. Contingencia: Mantener perfil sin habilitar, completar recuperación en ventana permitida y reforzar acompañamiento propio. Plazo: Antes de habilitar la persona y fuera del peak protegido; preparación anual en noviembre.

Se exige la evidencia siguiente: Asistencia, ejercicio aprobado, recuperación y cuadrante de acompañamiento. Paquetes vinculados: 1.11.2/3/4/6/7; 1.12.6/10. Reserva: Cohortes y estabilización de SD7; no convertir asistencia en certificación. Fundamento: C07 13.3, pp. 29–30; BT cap. 22, p. 38; SD7 formación y adopción. Estado: Identificado; tratamiento comprometido; residual objetivo no verificado. \parencites[numeral 13.3, pp.~29--30]{sd8-c07}[cap.~22, p.~38]{sd8-bt}

\subsection{R-15: Pilotos o evidencia ambiental fuera de ventana}
Dependencias de las innovaciones y ciclos naturales; medición tardía. No se forma la serie o no se evalúa el piloto antes de su fecha comprometida. Se demora una innovación o queda insuficiente evidencia para la temporada 2029. La evaluación inicial es $P=3$, $I=4$ y $E=12$; el objetivo posterior al tratamiento es $P_o=2$, $I_o=4$ y $E_o=8$, sin acreditación de residual observado.

Ignacio Silva (Proyecto) responde por el riesgo. Disparador: Serie incompleta, piloto fuera del ciclo natural o dependencia de innovación sin recepción. Mantener fechas y dependencias de SD13, habilitar medición temprana y verificar datos/pilotos con criterios propios. Contingencia: Reforzar captura y evaluación en ventana válida; separar prototipo y materialización contractual y escalar sin inventar resultados. Plazo: Antes del piloto/ciclo de SD13; seguimiento mensual de serie previa a temporada 2029.

Se exige la evidencia siguiente: Serie íntegra, criterios, calendario y acta del piloto aplicable. Paquetes vinculados: 1.9; ventanas de SD13; medición y evidencia E1/E2. Reserva: Ventanas naturales no son HH ni holgura disponible; autores y responsables de SD13 se conservan. Fundamento: C07 13.2, p. 29; BA arts. 28–30, pp. 19–20; SD13/SD7 innovaciones. Estado: Identificado; tratamiento comprometido; residual objetivo no verificado. \parencites[numeral 13.2, p.~29]{sd8-c07}[arts.~28--30, pp.~19--20]{sd8-ba}

\subsection{R-16: Fuentes, licencias o custodia incompletas}
Depósito por liberación omitido, titularidad incorrecta o custodia sin habilitar. El CLIENTE no recibe el paquete reproducible o no puede activar su custodia. Se incumplen propiedad, continuidad y transferencia de salida. La evaluación inicial es $P=3$, $I=4$ y $E=12$; el objetivo posterior al tratamiento es $P_o=2$, $I_o=3$ y $E_o=6$, sin acreditación de residual observado.

Ignacio Silva (Proyecto) responde por el riesgo. Disparador: Manifiesto/fuentes ausentes, licencia no compatible o comprobante semestral faltante. Repositorio del CLIENTE, verificación de depósito en cada liberación y contrato de custodia independiente con condiciones de liberación. Contingencia: Bloquear promoción incompleta, reconstruir paquete desde commit verificado y subsanar custodia/licencia antes de declarar recepción. Plazo: Cada liberación; custodia inicial y actualización semestral; revisión antes de salida.

Se exige la evidencia siguiente: Titularidad, acceso, fuentes/digest, licencia y comprobante del custodio. Paquetes vinculados: 1.2.1.1; 1.1.3.2/3; 1.12.11; 1.12.8.3. Reserva: Custodia de SD7 y puerta de liberación T-10; sin doble depósito presupuestado como nueva construcción. Fundamento: BA art. 84, p. 43; SD6 T-10; SD7 custodia. Estado: Identificado; tratamiento comprometido; residual objetivo no verificado. \parencite[art.~84, p.~43]{sd8-ba}

\section{Consolidación de los riesgos de IN1--IN5}\label{sec:t16-innovaciones}
Este complemento conserva R-01 a R-16 y reúne dieciocho mecanismos específicos, con los códigos R-IN1 a R-IN5. Cada mecanismo se vincula a un riesgo padre sin sustituir su titular ni valoración; su causa, disparador y acción se revisan por separado. La fuente versionada es \path{componentes/registro_riesgos_innovaciones.csv}. La exposición del padre no se recalcula sumando productos ordinales. Si varios mecanismos materializan el mismo evento, Proyecto conserva una incidencia, una imputación y un consumo de reserva.

Se aplica la escala de SD8: probabilidad alta/media/baja preliminar corresponde a 4/3/2; impacto medio/alto/muy alto, a 3/4/5. Esta correspondencia sirve para la evaluación de oferta, no para asignar frecuencias. El impacto toma el mayor efecto: R-IN2-03 se eleva a 5 por seguridad de personas y R-IN3-03/R-IN4-01/02/R-IN5-03 por privacidad o aceptación contractual, según la causa desarrollada. Los objetivos residuales son metas; el residual observado, su fecha y evaluador competente se incorporarán sólo con evidencia de eficacia.

Todos los mecanismos tienen estado identificado y tratamiento comprometido, sin pruebas ni aceptación ya ejecutadas. Sus plazos son hitos futuros del calendario vigente de SD13/SD7. Proyecto revisará cada ficha en todos los Comités de Proyecto y al disparador, registrando versión, fecha real de revisión, exposición, acción siguiente y aceptación competente. Las condiciones impeditivas bloquean la función o el reconocimiento económico; una puntuación menor no autoriza una excepción obligatoria. La identificación de estos riesgos cumple el alcance general de RT-19.04 y no declara cumplidos los criterios de beneficio o madurez de RT-26.

Las acciones usan paquetes existentes y se nivelan por persona en T-15 antes de seleccionar trabajo. No se crea una reserva de dinero ni se imputa de nuevo un ahorro de IN1 a IN4. Una respuesta que requiera esfuerzo adicional se dimensiona y financia según SD8 y OE; sólo CC-6 autoriza un cambio contractual. La retirada de una extensión conserva las prestaciones obligatorias.
\subsection{R-IN1-01: Actualización insuficiente por contratistas}
Enlace T-16: R-14. Causa: Se conserva la coordinación telefónica sin ingresar sus cambios. Evento: Actualización insuficiente por contratistas. Consecuencia: El historial queda incompleto y el armador recibe un estado obsoleto.

Evaluación inicial: $P=4$, $I=4$ y $E=16$. Objetivo residual: $P_o=3$, $I_o=3$ y $E_o=9$, sin residual observado acreditado. Titular: José Lara Arce (Implantación).

Disparador: Cambio comunicado sin registro durante más de un día hábil. Tratamiento: Formar contratistas y marina; comprobar tareas y recordatorios con fuente de cada cambio. Contingencia: Activar registro asistido de la marina, conservando autor, fuente e instante del hecho; conciliar con el contratista.

Hito o plazo: Formación al término del mes 18; control diario del piloto 19 y seguimiento 21--56. Evidencia exigida: Comunicaciones, cambios pendientes, registro asistido y comprobación de concordancia. Paquetes: IN1.1.1; IN1.3.3; IN1.5.1. Estado: Identificado; tratamiento comprometido; residual objetivo no verificado. Reserva y capacidad: Acciones en paquete vinculado y T-15; sin HH libres ni reserva económica adicional; costo y refuerzo se concilian en OE con autorización. Para R-IN2-03 rige R-01, con consumo único por campaña.

\subsection{R-IN1-02: Participación insuficiente de armadores}
Enlace T-16: R-15. Causa: Adhesión voluntaria y pocas solicitudes elegibles. Evento: Participación insuficiente de armadores. Consecuencia: No se forma una comparación evaluable del beneficio.

Evaluación inicial: $P=3$, $I=3$ y $E=9$. Objetivo residual: $P_o=2$, $I_o=3$ y $E_o=6$, sin residual observado acreditado. Titular: José Lara Arce (Implantación).

Disparador: Menos de quince solicitudes elegibles al terminar la primera quincena del piloto o menos de treinta al cierre. Tratamiento: Orientar e ingresar solicitudes de forma asistida sin crear demanda artificial; conservar excluidas y no respondidas. Contingencia: Reforzar convocatoria; declarar beneficio no demostrado si falta muestra; cualquier extensión seguirá CC-6 y ventana autorizada.

Hito o plazo: Primera quincena y cierre del piloto del mes 19. Evidencia exigida: Listado completo, elegibilidad, motivos de exclusión y acta de resultado, incluidos resultados nulos. Paquetes: IN1.5.1; IN1.5.4; IN1.5.6. Estado: Identificado; tratamiento comprometido; residual objetivo no verificado. Reserva y capacidad: Acciones en paquete vinculado y T-15; sin HH libres ni reserva económica adicional; costo y refuerzo se concilian en OE con autorización. Para R-IN2-03 rige R-01, con consumo único por campaña.

\subsection{R-IN1-03: Integración incompatible de solicitudes y órdenes}
Enlace T-16: R-03. Causa: Contratos, identidades o versiones difieren entre IN1, CTR y VAR. Evento: Integración incompatible de solicitudes y órdenes. Consecuencia: Solicitud no vinculada, orden duplicada o transición incoherente.

Evaluación inicial: $P=3$, $I=4$ y $E=12$. Objetivo residual: $P_o=2$, $I_o=4$ y $E_o=8$, sin residual observado acreditado. Titular: Fernando Guerrero (Arquitectura).

Disparador: Interfaz crítica no supera pruebas al cierre del mes 16; duplicado o rechazo de versión no controlado. Tratamiento: Probar contratos, idempotencia, concurrencia, habilitación y recuperación de avisos; validar vínculo con orden autorizada. Contingencia: Bloquear flujo adicional afectado; conservar solución base y diario, corregir y repetir pruebas; gestionar ajuste de piloto mediante CC-6.

Hito o plazo: Aceptación técnica en mes 16; comprobación previa a piloto 19 y materialización 21. Evidencia exigida: Versiones de contrato, casos negativos, historial y conciliación sin órdenes duplicadas. Paquetes: IN1.1.2; IN1.3.1--3; IN1.4.1. Estado: Identificado; tratamiento comprometido; residual objetivo no verificado. Reserva y capacidad: Acciones en paquete vinculado y T-15; sin HH libres ni reserva económica adicional; costo y refuerzo se concilian en OE con autorización. Para R-IN2-03 rige R-01, con consumo único por campaña.

\subsection{R-IN1-04: Acceso a solicitudes de otra embarcación}
Enlace T-16: R-10. Causa: Autorización insuficiente de armador o contratista asignado. Evento: Acceso a solicitudes de otra embarcación. Consecuencia: Divulgación o modificación ajena a la finalidad autorizada.

Evaluación inicial: $P=2$, $I=5$ y $E=10$. Objetivo residual: $P_o=1$, $I_o=5$ y $E_o=5$, sin residual observado acreditado. Titular: Bruno Toro (Seguridad).

Disparador: Cualquier acceso indebido confirmado o prueba negativa fallida. Tratamiento: Validar vínculo vigente y ámbito en servidor; ensayar separación entre armadores y contratistas con auditoría. Contingencia: Suspender función afectada, conservar evidencia y activar incidente; revisar alcance, corregir y repetir pruebas antes de reactivar.

Hito o plazo: Antes del piloto y de cada liberación; contención inmediata. Evidencia exigida: Matriz de roles, pruebas positivas/negativas, auditoría e informe de reexamen. Paquetes: IN1.3.2; IN1.4.2. Estado: Identificado; tratamiento comprometido; residual objetivo no verificado. Reserva y capacidad: Acciones en paquete vinculado y T-15; sin HH libres ni reserva económica adicional; costo y refuerzo se concilian en OE con autorización. Para R-IN2-03 rige R-01, con consumo único por campaña.

\subsection{R-IN2-01: Cohorte de invernada insuficiente}
Enlace T-16: R-15. Causa: Adhesión voluntaria o convocatoria sin grupo estable. Evento: Cohorte de invernada insuficiente. Consecuencia: No se demuestra coordinación preventiva ni beneficio estacional.

Evaluación inicial: $P=3$, $I=4$ y $E=12$. Objetivo residual: $P_o=2$, $I_o=4$ y $E_o=8$, sin residual observado acreditado. Titular: José Lara Arce (Implantación).

Disparador: Al cierre de inscripción del mes 30 no se reúnen las veinte embarcaciones planificadas con plan y consentimiento. Tratamiento: Convocar y formar armadores; confirmar planes y disponibilidad sin imponer contratista ni trabajos. Contingencia: Proponer convocatoria adicional o alcance evaluable mediante CC-6; documentar limitación y no declarar alcanzada la evaluación de veinte barcos.

Hito o plazo: Inscripción y planes en mes 30, antes de invernada del 31. Evidencia exigida: Cohorte, consentimientos, planes, bajas y acta de decisión sobre evaluación. Paquetes: IN2.1.2; IN2.5.1. Estado: Identificado; tratamiento comprometido; residual objetivo no verificado. Reserva y capacidad: Acciones en paquete vinculado y T-15; sin HH libres ni reserva económica adicional; costo y refuerzo se concilian en OE con autorización. Para R-IN2-03 rige R-01, con consumo único por campaña.

\subsection{R-IN2-02: Atraso del contratista respecto de botadura}
Enlace T-16: R-15. Causa: Trabajo o disponibilidad no confirmado antes de la ventana. Evento: Atraso del contratista respecto de botadura. Consecuencia: Se llega a botadura con mantenimiento pendiente.

Evaluación inicial: $P=4$, $I=4$ y $E=16$. Objetivo residual: $P_o=3$, $I_o=3$ y $E_o=9$, sin residual observado acreditado. Titular: José Lara Arce (Implantación).

Disparador: Revisión a tres semanas de cada botadura detecta trabajo pendiente sin fecha viable. Tratamiento: Coordinar líder funcional y administración de varadero; confirmar disponibilidad y los cinco lotes de cuatro embarcaciones. Contingencia: Proponer reprogramación o sustitución del contratista con aceptación del armador; preservar habilitación y seguridad, sin prometer trabajos ajenos ejecutados.

Hito o plazo: Cinco revisiones al menos tres semanas antes de botaduras septiembre--noviembre de 2029. Evidencia exigida: Confirmaciones, avance, acuerdos del armador y calendario de pendientes por barco. Paquetes: IN2.3.1; IN2.5.2; IN2.5.3. Estado: Identificado; tratamiento comprometido; residual objetivo no verificado. Reserva y capacidad: Acciones en paquete vinculado y T-15; sin HH libres ni reserva económica adicional; costo y refuerzo se concilian en OE con autorización. Para R-IN2-03 rige R-01, con consumo único por campaña.

\subsection{R-IN2-03: Suspensión meteorológica durante invernada}
Enlace T-16: R-01. Causa: Aviso marítimo o condición insegura coincidente con trabajos programados. Evento: Suspensión meteorológica durante invernada. Consecuencia: Suspensión obligatoria y pérdida de ventana; exposición de personas si se continúa.

Evaluación inicial: $P=4$, $I=5$ y $E=20$. Objetivo residual: $P_o=4$, $I_o=3$ y $E_o=12$, sin residual observado acreditado. Titular: Ignacio Silva (Proyecto).

Disparador: Aviso oficial de marejadas o condición insegura; reserva disponible insuficiente para reprogramar. Tratamiento: Integrar ventana y aviso con R-01; comprobar reserva por campaña y red de todos los frentes antes de seleccionar trabajo. Contingencia: Suspender todo trabajo programado; consumir reserva una sola vez; decidir refuerzo y pendientes tras cese y autorización, sin extender piloto automáticamente.

Hito o plazo: Control previo a cada campaña y revisión inmediata al aviso en meses 31--36. Evidencia exigida: Aviso, bloqueo, saldo de reserva, agenda revisada y autorización de reanudación. Paquetes: IN2.5.2; IN2.5.3; PL-R01. Estado: Identificado; tratamiento comprometido; residual objetivo no verificado. Reserva y capacidad: Acciones en paquete vinculado y T-15; sin HH libres ni reserva económica adicional; costo y refuerzo se concilian en OE con autorización. Para R-IN2-03 rige R-01, con consumo único por campaña.

\subsection{R-IN2-04: IN1 o interfaces no disponibles para IN2}
Enlace T-16: R-03. Causa: La coordinación preventiva depende de un servicio o contrato aún no recibido. Evento: IN1 o interfaces no disponibles para IN2. Consecuencia: No se enlazan planes, solicitudes o avance de la invernada.

Evaluación inicial: $P=3$, $I=4$ y $E=12$. Objetivo residual: $P_o=2$, $I_o=4$ y $E_o=8$, sin residual observado acreditado. Titular: Fernando Guerrero (Arquitectura).

Disparador: Prueba IN2/IN1 fallida o dependencia no recibida antes de iniciar invernada. Tratamiento: Verificar contrato, estados y disponibilidad de IN1 antes de habilitar el plan preventivo. Contingencia: Mantener coordinación controlada por marina y solución base; conciliar referencias; bloquear extensión no verificada hasta repetir prueba.

Hito o plazo: Antes del piloto estacional del mes 31 y ante cambios de IN1. Evidencia exigida: Contrato recibido, prueba de enlace y conciliación del registro alternativo. Paquetes: IN2.3.2; IN2.4.1. Estado: Identificado; tratamiento comprometido; residual objetivo no verificado. Reserva y capacidad: Acciones en paquete vinculado y T-15; sin HH libres ni reserva económica adicional; costo y refuerzo se concilian en OE con autorización. Para R-IN2-03 rige R-01, con consumo único por campaña.

\subsection{R-IN3-01: Campos OCR incorrectos por formato o calidad}
Enlace T-16: R-03. Causa: Documento ilegible o formato cambiado; extractor propone valores erróneos. Evento: Campos OCR incorrectos por formato o calidad. Consecuencia: Póliza queda asociada con vigencia, monto o embarcación incorrectos.

Evaluación inicial: $P=4$, $I=4$ y $E=16$. Objetivo residual: $P_o=3$, $I_o=3$ y $E_o=9$, sin residual observado acreditado. Titular: Fernando Guerrero (Arquitectura; coordinación con Líder de Datos).

Disparador: Error crítico de campo o recurrencia por formato en prueba o piloto. Tratamiento: Líder de Datos aplica reglas, fragmentos de origen y evaluación por formato; conservar original y propuesta separada del dato aprobado. Contingencia: Suspender extracción del formato afectado y usar ingreso manual revisado; corregir y repetir pruebas antes de rehabilitar.

Hito o plazo: Pruebas del mes 16; evaluación intercalada manual/asistida en meses 18--19 y cada cambio de formato. Evidencia exigida: Original, propuesta, corrección, validación de campos y métricas desagregadas por formato. Paquetes: IN3.2.2; IN3.4.1. Estado: Identificado; tratamiento comprometido; residual objetivo no verificado. Reserva y capacidad: Acciones en paquete vinculado y T-15; sin HH libres ni reserva económica adicional; costo y refuerzo se concilian en OE con autorización. Para R-IN2-03 rige R-01, con consumo único por campaña.

\subsection{R-IN3-02: Revisión humana insuficiente de extracción}
Enlace T-16: R-14. Causa: Sesgo de automatización, formación incompleta o revisión apresurada. Evento: Revisión humana insuficiente de extracción. Consecuencia: Se aprueban datos incorrectos y se toman decisiones con antecedentes defectuosos.

Evaluación inicial: $P=3$, $I=4$ y $E=12$. Objetivo residual: $P_o=2$, $I_o=4$ y $E_o=8$, sin residual observado acreditado. Titular: José Lara Arce (Implantación; con verificación de Calidad).

Disparador: Un error crítico supera la revisión humana o revisor no acredita competencia de tarea. Tratamiento: Capacitar revisores con casos erróneos; exigir confirmación de campos y comprobar registros aprobados por QA. Contingencia: Pausar flujo afectado, revisar registros potencialmente comprometidos y reforzar revisión antes de reanudar; mantener ingreso manual.

Hito o plazo: Antes de habilitar revisores y durante piloto 18--19 y operación 21--56. Evidencia exigida: Ejercicio individual, revisión de campos, errores aprobados y reexamen de registros afectados. Paquetes: IN3.5.1; IN3.5.3. Estado: Identificado; tratamiento comprometido; residual objetivo no verificado. Reserva y capacidad: Acciones en paquete vinculado y T-15; sin HH libres ni reserva económica adicional; costo y refuerzo se concilian en OE con autorización. Para R-IN2-03 rige R-01, con consumo único por campaña.

\subsection{R-IN3-03: Tratamiento documental no autorizado}
Enlace T-16: R-10. Causa: Proveedor, ubicación, retención o permiso difieren de condiciones aprobadas. Evento: Tratamiento documental no autorizado. Consecuencia: Exposición de pólizas o tratamiento fuera de finalidad autorizada.

Evaluación inicial: $P=3$, $I=5$ y $E=15$. Objetivo residual: $P_o=2$, $I_o=5$ y $E_o=10$, sin residual observado acreditado. Titular: Bruno Toro (Seguridad).

Disparador: Condición de proveedor incompatible, permiso fallido o incidente documental. Tratamiento: Comprobar ubicación, retención, permisos y condiciones de procesamiento antes de enviar; conservar trazabilidad y acceso mínimo. Contingencia: Bloquear envío, activar respuesta de incidente y utilizar ingreso manual; rehabilitar sólo con condiciones y controles conformes.

Hito o plazo: Antes de contratar/habilitar procesamiento y ante cambios del proveedor o liberación. Evidencia exigida: Condiciones aprobadas, configuración, accesos positivos/negativos y evidencia de retención/borrado aplicable. Paquetes: IN3.1.2; IN3.4.2. Estado: Identificado; tratamiento comprometido; residual objetivo no verificado. Reserva y capacidad: Acciones en paquete vinculado y T-15; sin HH libres ni reserva económica adicional; costo y refuerzo se concilian en OE con autorización. Para R-IN2-03 rige R-01, con consumo único por campaña.

\subsection{R-IN3-04: Extractor o conectividad interrumpidos}
Enlace T-16: R-02. Causa: Dependencia remota no disponible durante ingreso. Evento: Extractor o conectividad interrumpidos. Consecuencia: Se acumulan pólizas pendientes o se duplica ingreso al recuperar.

Evaluación inicial: $P=3$, $I=3$ y $E=9$. Objetivo residual: $P_o=2$, $I_o=3$ y $E_o=6$, sin residual observado acreditado. Titular: José Lara Arce (Operación/SRE).

Disparador: Fallo del extractor, agotamiento de reintentos o enlace indisponible. Tratamiento: Supervisar fallos y ensayar alternativa manual, reintentos e idempotencia; separar pendientes de antecedentes aprobados. Contingencia: Ingresar manualmente con original y fuente; al recuperar, conciliar resultados pendientes sin duplicar ni sobrescribir aprobados.

Hito o plazo: Prueba previa a piloto y operación; actuación al fallo. Evidencia exigida: Registro de fallo, cola pendiente, ingreso manual y conciliación tras recuperación. Paquetes: IN3.3.3; IN3.4.1. Estado: Identificado; tratamiento comprometido; residual objetivo no verificado. Reserva y capacidad: Acciones en paquete vinculado y T-15; sin HH libres ni reserva económica adicional; costo y refuerzo se concilian en OE con autorización. Para R-IN2-03 rige R-01, con consumo único por campaña.

\subsection{R-IN3-05: Beneficio de extracción insuficiente}
Enlace T-16: R-15. Causa: Volumen bajo o revisión consume el ahorro de captura. Evento: Beneficio de extracción insuficiente. Consecuencia: Costos incrementales superan el beneficio justificable.

Evaluación inicial: $P=3$, $I=3$ y $E=9$. Objetivo residual: $P_o=2$, $I_o=3$ y $E_o=6$, sin residual observado acreditado. Titular: Ignacio Silva (Proyecto).

Disparador: Muestra de 120 pólizas distintas incompleta o tiempo/volumen observado no sustenta beneficio neto. Tratamiento: Medir manual/asistido con revisión incluida, por formato; registrar esfuerzo y contrastar costos incrementales exclusivamente en OE. Contingencia: Declarar beneficio no demostrado y mantener ingreso manual; ajustar o deshabilitar extensión sin ampliar alcance automáticamente.

Hito o plazo: Cierre del piloto 18--19 y revisión previa a materialización 21. Evidencia exigida: Muestra y tiempos con revisión, resultados adversos y referencia a evaluación económica autorizada. Paquetes: IN3.1.1; IN3.5.3. Estado: Identificado; tratamiento comprometido; residual objetivo no verificado. Reserva y capacidad: Acciones en paquete vinculado y T-15; sin HH libres ni reserva económica adicional; costo y refuerzo se concilian en OE con autorización. Para R-IN2-03 rige R-01, con consumo único por campaña.

\subsection{R-IN4-01: Atribución inválida o pago por comparación no equivalente}
Enlace T-16: R-15. Causa: Mezcla, dotación o mejoras ajenas cambian respecto de la base; medición no aprobada. Evento: Atribución inválida o pago por comparación no equivalente. Consecuencia: Reconocimiento variable indebido o desacuerdo contractual.

Evaluación inicial: $P=3$, $I=5$ y $E=15$. Objetivo residual: $P_o=2$, $I_o=5$ y $E_o=10$, sin residual observado acreditado. Titular: Ignacio Silva (Proyecto).

Disparador: Base cero, menos de treinta solicitudes por condición, datos no comparables o discrepancia de cálculo/atribución. Tratamiento: Fijar ficha y composición antes de medir; conservar intervalos activos, fuentes, exclusiones y cálculo; separar preparación Synaptix de verificación cliente. Contingencia: Asignar factor cero hasta disponer de evidencia aceptada; corregir prospectivamente y repetir evaluación autorizada; mantener servicio fijo y SLA.

Hito o plazo: Ficha antes de base del mes 18; evaluación 19; conciliación mensual 21--56. Evidencia exigida: Ficha aprobada, tiempos por persona, cálculos versionados, conciliación y decisión del cliente; valorización sólo OE. Paquetes: IN4.1; IN4.3; IN4.4. Estado: Identificado; tratamiento comprometido; residual objetivo no verificado. Reserva y capacidad: Acciones en paquete vinculado y T-15; sin HH libres ni reserva económica adicional; costo y refuerzo se concilian en OE con autorización. Para R-IN2-03 rige R-01, con consumo único por campaña.

\subsection{R-IN4-02: Omisión de casos difíciles o cierre apresurado por incentivo}
Enlace T-16: R-15. Causa: Selección de tareas fáciles o presión para aparentar menor esfuerzo. Evento: Omisión de casos difíciles o cierre apresurado por incentivo. Consecuencia: Se ocultan errores/abiertas, empeora el servicio y se reclama incentivo sin resultado válido.

Evaluación inicial: $P=3$, $I=5$ y $E=15$. Objetivo residual: $P_o=2$, $I_o=5$ y $E_o=10$, sin residual observado acreditado. Titular: Ignacio Silva (Proyecto; con verificación de Calidad).

Disparador: Solicitud elegible ausente sin motivo, cierre no sustentado, aumento de errores/reaperturas o caída de proporción de cierre. Tratamiento: Calidad contrasta universo elegible con AUT/VAR/CTR, incluidas abiertas y excluidas; fija complejidad y plazo antes de resultados; compara calidad por condición. Contingencia: Suspender reconocimiento y aplicar factor cero; revisar universo y cierres, corregir y repetir evaluación sin borrar versión; sostener operación y obligaciones de atención.

Hito o plazo: Antes de cerrar piloto 19 y cada conciliación mensual 21--56. Evidencia exigida: Universo conciliado, motivos de exclusión, abiertas, errores, reaperturas y control de cierre por condición. Paquetes: IN4.2; IN4.3; IN4.4. Estado: Identificado; tratamiento comprometido; residual objetivo no verificado. Reserva y capacidad: Acciones en paquete vinculado y T-15; sin HH libres ni reserva económica adicional; costo y refuerzo se concilian en OE con autorización. Para R-IN2-03 rige R-01, con consumo único por campaña.

\subsection{R-IN5-01: Reducción ambiental aparente por datos o selección sesgados}
Enlace T-16: R-15. Causa: Adhesión voluntaria, lecturas faltantes o cambios de ocupación, clima y lavado. Evento: Reducción ambiental aparente por datos o selección sesgados. Consecuencia: Se publica ahorro no atribuible al plan o sustentado en consumo/permanencia incorrectos.

Evaluación inicial: $P=3$, $I=4$ y $E=12$. Objetivo residual: $P_o=2$, $I_o=4$ y $E_o=8$, sin residual observado acreditado. Titular: Ignacio Silva (Proyecto; con verificación de Calidad).

Disparador: Lectura/permanencia no confiable, base o denominador cero, composición distinta o período sin evidencia. Tratamiento: Comparar mismas embarcaciones, conservar asociación histórica de medidor y permanencia validada; fijar criterios y registrar faenas, clima y resultados adversos. Contingencia: No publicar reducción; declarar período no evaluable y corregir/capturar en ventana válida sin imputar cero ni atribuir causalidad por simple diferencia.

Hito o plazo: Base del mes 18, piloto 19; antes de publicación y preparación de evidencia ambiental 2029. Evidencia exigida: Series conciliadas, asociaciones históricas, permanencia, exclusiones, contexto y cálculo reproducible. Paquetes: IN5.1; IN5.3; IN5.4. Estado: Identificado; tratamiento comprometido; residual objetivo no verificado. Reserva y capacidad: Acciones en paquete vinculado y T-15; sin HH libres ni reserva económica adicional; costo y refuerzo se concilian en OE con autorización. Para R-IN2-03 rige R-01, con consumo único por campaña.

\subsection{R-IN5-02: Participación o continuidad insuficientes en planes ambientales}
Enlace T-16: R-14. Causa: Plan poco comprensible o abandono voluntario de armador/charter. Evento: Participación o continuidad insuficientes en planes ambientales. Consecuencia: No se ejecutan acciones ni se evalúa utilidad del plan.

Evaluación inicial: $P=3$, $I=3$ y $E=9$. Objetivo residual: $P_o=2$, $I_o=3$ y $E_o=6$, sin residual observado acreditado. Titular: José Lara Arce (Implantación).

Disparador: Cohorte no evaluable, retiro o acción sin responsable/plazo/evidencia. Tratamiento: Validar acciones breves con armadores y charter; explicar responsabilidades y retiro sin pérdida de servicio básico. Contingencia: Ajustar o desactivar plan adicional y registrar resultado no demostrado; preservar registros obligatorios y cierre de no conformidades.

Hito o plazo: Antes de acciones del piloto 19 y en seguimiento 21--56. Evidencia exigida: Adhesiones/retiros, comprensión de tarea, acciones, plazos y evaluación de utilidad. Paquetes: IN5.1; IN5.3; IN5.4. Estado: Identificado; tratamiento comprometido; residual objetivo no verificado. Reserva y capacidad: Acciones en paquete vinculado y T-15; sin HH libres ni reserva económica adicional; costo y refuerzo se concilian en OE con autorización. Para R-IN2-03 rige R-01, con consumo único por campaña.

\subsection{R-IN5-03: Acceso entre flotas o uso fuera de finalidad ambiental}
Enlace T-16: R-10. Causa: Vínculo histórico/consentimiento o autorización por embarcación no se comprueba. Evento: Acceso entre flotas o uso fuera de finalidad ambiental. Consecuencia: Se expone historial de otra flota o se habilita finalidad no autorizada.

Evaluación inicial: $P=2$, $I=5$ y $E=10$. Objetivo residual: $P_o=1$, $I_o=5$ y $E_o=5$, sin residual observado acreditado. Titular: Bruno Toro (Seguridad).

Disparador: Prueba de aislamiento fallida, consentimiento modificado no aplicado o acceso indebido. Tratamiento: Verificar vínculo vigente, finalidad y autorización por embarcación; ensayar flotas separadas y propagación de restricciones con auditoría. Contingencia: Suspender vista afectada, conservar evidencia y activar incidente; corregir acceso y repetir aislamiento antes de rehabilitar.

Hito o plazo: Antes del piloto, cada liberación y al cambio de vínculo/consentimiento. Evidencia exigida: Casos de separación entre flotas, auditoría de consultas/cambios y prueba de restricción propagada. Paquetes: IN5.2; IN5.3; IN5.4. Estado: Identificado; tratamiento comprometido; residual objetivo no verificado. Reserva y capacidad: Acciones en paquete vinculado y T-15; sin HH libres ni reserva económica adicional; costo y refuerzo se concilian en OE con autorización. Para R-IN2-03 rige R-01, con consumo único por campaña.


\section{Revisión, evidencia y aceptación}
Proyecto revisará este registro en cada Comité de Proyecto y ante cambios, incidentes o ventanas críticas. El titular aporta resultado y residual observado; Calidad verifica eficacia; la autoridad competente decide aceptación dentro de sus atribuciones. Seguridad y SRE conservan el procedimiento de incidente y la Contraparte conserva recepción. Las actas se emiten en dos días hábiles y desvíos mayores a 10\,\% de un hito se comunican en cinco días hábiles con recuperación. \parencites[art.~71, pp.~37--38]{sd8-ba}[RT-19.04/09/10, pp.~33--34]{sd8-bt}
El proceso aplica ISO 31000:2018 y utiliza análisis de modos de falla y sensibilidad de SD8, con IEC 31010 como referencia para seleccionar técnicas. Las relaciones, escalas y decisiones se mantienen versionadas, sin atribuir una certificación ni evaluación humana ya realizada. \parencite{sd8-iso31000,sd8-31010}
\CierreSubdocumento
\end{document}
